\chapter{36}



Joh sass mit einem Kugelschreiber in der Hand über Papiere gebeugt am
langen Tisch. Die Mittagswärme war bereits in die Küche unterwegs. Wie
sich Louis neben Joh setzte, verflog seine Unsicherheit. Joh schaute freundlich auf und blickte ihm in die Augen.
Dann schob er ihm eine Obstschale zu und schenkte ihm Tee ein.
Abschliessend legte er frische Pfefferminzblätter obendrauf. Louis
schaute ihm dabei begeistert zu. 

«Danke, Joh,» Louis nahm sich eine Birne aus der Schale, «die kommt vom
Birnbaum draussen, ja?»

Joh nickte und zog die Augenbrauen hoch.

«Gut kombiniert, Eingangsprüfung bestanden.»

Joh wirkte vergnügt. Schliesslich setzte er einen geschäftigen Blick
auf.

«Nun, Ciro, ohne grosse Umschweife. Du überzeugst uns. Manuela und ich
sind uns einig. Wie haben uns gestern besprochen. Wenn du willst, bieten
wir dir diese temporäre Anstellung an. Was meinst du?»

«Das freut mich. Und wenn dann noch was zur Deckung meiner laufenden Kosten bleibt - dann, dann bin ich gerne dabei.»

«Klar, versteht sich von selbst,» er grinste, «als obdachlosen
Streuner schätze ich dich nicht ein.»

«Dann hab ich ja nochmal Glück gehabt.»

Louis grinste mager zurück.

Die folgende Stunde gehörte Louis’ Arbeitsvertrag. Joh führte ihn Punkt für Punkt durch das Dokument.
Nichts schien er auszulassen. Neben Sozial- und Unfallversicherung,
Krankenkasse und sogar alpgerechter Kleidung, wurden auch die
Arbeitsanforderungen präzise definiert. Als kurz in der Schweiz
Arbeitender, als Saisonnier, sei einiges zu klären. Joh schien konzentriert nachzudenken.

«Bei der Anstellung von Saisonniers mit ausländischem Wohnsitz habe ich bisher keine Erfahrung, hm, das braucht eine Anmeldung an deinem Wohnsitz in Italien,» Louis zuckte zusammen, versuchte sich aber seinen Schrecken nicht anmerken zu lassen, «Moment,» Joh drehte gedankenversunken den Kugelschreiber hin und her, «also, ich mache dir einen Vorschlag. Meine Arbeitsbelastung ist gerade gross, aber bis Mitte September sollte ich deine Anmeldung und die restlichen Formalitäten erledigt haben. Wenn es für dich in Ordnung ist, starten wir unter den vereinbarten Bedingungen und ich stelle dich rückwirkend an. Also, keine Sorge, ich bezahle dich später natürlich vom Tag deines Einsatzes an. Was meinst du?»

Was sollte er denn meinen? Es war wohl das Beste, das ihm gerade passieren konnte. Jeder Tag ohne Anmeldung verlängerte seine Gnadenfrist.

«Bin gerne damit einverstanden, Joh. Nur, gerade bin ich knapp bei Kasse, aber - ,» Joh unterbrach ihn sanft, «natürlich brauchst du die nächsten Tage bei uns kein Geld, Ciro, aber dazu komme ich gleich, ok?»

Louis nickte schnell.

Ein weiterer Abschnitt war alleine dem Verhaltenskodex für den Bereich \emph{Alp und Umgang mit den Tieren}
gewidmet. Louis schien, als ob Joh Louis’ Reaktion nun besonders
aufmerksam studierte.
Er las parallel zu Johs Ausführungen auf einem Doppel mit. Immer
wieder schielte er zum Ende des Papiers. Er suchte eine Zahl, die er
nicht fand. Was sollte er tun, wenn kaum etwas raussprang? Und, überhaupt, wieviele Tage blieben ihm als Ciro auf diesem Hof? Joh schien seine Ungeduld aufzufallen.

«Bald haben wir’s, Ciro,» er legte seine Hand auf Louis’ Arm, «ein gut
abgesprochener Vertrag ist uns wichtig.»

Louis zog die Mundwinkel hoch. Von draussen erreichte sie ein Geräusch, als ziehe jemand feinstes
Schleifpapier über die Zufahrt. Mit Blick aus dem Fenster entdeckte
Louis ein Fahrzeug, das sich der \emph{Sunnewaid} näherte. Es stoppte
direkt vor dem Haus.
Eine weibliche Stimme ertönte.

«Hallihallo, Lieferung.»

«Bin gleich zurück, Ciro.»

Joh verschwand vor das Haus. Der Unterbruch schürte Louis’ Ungeduld zusätzlich und sein rechter Fuss begann nervös auf- und ab zu wippen. Nach einer Weile kam Joh mit einer sperrigen
Kiste auf den Armen zurück und sass Louis bald wieder gegenüber am Tisch.

«Wir richten uns nach den Empfehlungen des \emph{Alpwirtschaftlichen
Verbands,} was unsere Bezahlung betrifft. In Abzug von Kost, Logie,
Krankenkasse, Sozialversicherung und Arbeitskleidung, kann ich dir monatlich drei-eins anbieten. Was meinst du, Deal?»

«Du meinst dreitausendeinhundert Franken?"

Louis schaute ihn ungläubig an. Joh schien verwirrt.

«Ja.»

«Gut, klar, das passt, ich danke dir.»

«Gerne. Nun denn, Arbeitsbeginn, Montag, 27. August bis vorläufig Mitte Dezember, einverstanden?»

Louis nickte einmal mehr, drückte Johs Hand und der trug schliesslich den vereinbarten Betrag
handschriftlich ein. Das waren erstaunliche Bedingungen. Ganz offensichtlich ein Schweizer Lohn 
und nicht vergleichbar mit seinem Gehalt in Italien. Sofort fiel ihm 
der hohe Bahnpreis seiner Reise ein. Anderes Land, andere Verhältnisse. 
Bis Mitte Dezember? Louis blieb nichts anderes als laufen zu lassen, was sich hier gerade abspielte. Auf jeden Fall aber musste er einen Weg finden, Marco über seine Auszeit zu informieren. Dabei hoffte er, dass Aldo Mathieu flexibel genug war, seinen Ausfall im Ristorante weiterhin zu überbrücken.

Nach einer ersten, grösseren Zahlung zu gegebener Zeit würde Louis anschliessend wöchentlich entlöhnt. Es sei bei ihnen gängige Praxis,
meinte Joh, da sie immer auch wieder temporäre Mitarbeiter für einzelne
Wochen anstellten. Jeweils am Samstag würde Louis der Lohn ausbezahlt
bekommen. In bar, wie Louis wünschte. Joh war kurz irritiert. Grundsätzlich gingen die Löhne seiner Mitarbeiter per Banküberweisung raus, hatte er erklärt. 
«Aber gut, einigen wir uns auf Barzahlung. Toll, Ciro, ich freue mich auf unsere Zusammenarbeit. Wir können loslegen. Vielleicht komme ich sogar früher dazu, mich um die Details zu kümmern, hoff ich jedenfalls,» er runzelte die Stirn, «nun brauche ich nur noch deinen Pass und lege die Kopie dem Vertrag bei, okay?»

«Natürlich.»

Louis stand auf. Steuern? All die Formalitäten? Was bedeutete das im Zusammenhang mit seiner falschen Identität? Er holte das rote Büchlein. Bis Mitte September wollte Joh die Formalitäten regeln? Nervös legte
er Joh das Dokument auf den Tisch. Wenn er Glück hatte, blieben ihm zwei Wochen, um wieder zu verschwinden. Der Gedanke war erschreckend. Joh stand pfeifend vor einem
Kopiergerät. Nach einem kurzen Blick in den Pass machte er die Kopien.
Er reichte Louis das Büchlein zurück. Joh vertraute ihm. Louis war fürs Erste 
erleichtert. Gleichzeitig schämte er sich angesichts der Fälschung. Er
beging klaren Betrug. 

Im Anschluss bat Joh Louis in den Vorraum des Stalls. Louis sollte sich
Arbeitskleidung und Schuhe aussuchen.
Neu ausgestattet stand er kurz darauf vor Joh. Er fühlte sich wie eine
Schaufensterpuppe, deren Ausstaffierung Joh prüfend klassierte. Als
hätten sie ein Zeichen vereinbart, lachten beide schallend.

«Du siehst aus wie der perfekte Regenschäfer der Alpen, toll,» Joh
grinste, «aber du wirst sehen, die Ausrüstung braucht’s.»

Louis schaute an sich hinunter und blieb fasziniert an den Bergschuhen
hängen. Sie hielten seine Füsse wie Schraubstöcke zusammen.

«Nichts gegen deine Trekkingschuhe, Ciro, aber solche da,» und Joh
zeigte auf Louis’ Füsse, «sind ein Muss. Der perfekte Schuh ist oben am
Gletscher das A und O. Der \emph{Sportiva} ist unser Favorit, Kategorie
C, fels- und gletschertauglich, stocksteife Sohle, steigeisenfest.»

Joh führte sich gerade wie ein Verkäufer auf. Die Grösse passte. Und als
Joh Louis’ zögerliches Herumgehen verfolgte, schien er ihn aufmuntern zu
wollen.

«Du wirst dich daran gewöhnen und den Schuh schätzen. Der war nur ein
paar Tage unterwegs, ist somit schon eingetragen und hat auf dich
gewartet.»

Joh zwinkerte Louis zu.

Schliesslich sassen sie nebeneinander. Draussen auf der Treppe zur
Küche. Beide eine Tasse Kaffee in der Hand. Fast schon enthusiastisch
erklärte Joh Louis den Ablauf des nächsten Tages.

«Um sieben steigen wir zum Gletscher hoch. Die beste Zeit, um vor der
grossen Hitze oben zu sein. Manuela kommt morgen noch mal mit. Zu dritt
schaffen wir es leichter, die Dinge hochzubringen, die du brauchst. Ich
werde dir Mazi vorstellen und er wird dich mit deinen Aufgaben vertraut
machen. Und ja, dann führe ich dich zur kleinen Hütte, die dir gehört,
wenn du oben bist. Mazi hat auf der andern Seite sein Alpheim,» dabei
schob Joh wieder die Brille hoch.

Louis’ Wunsch, seine Gitarre mit hochzunehmen, liess Joh im ersten
Moment verstummen..

«Oh, du hast Grosses vor,» er hob die Augenbrauen, «wenn du abends
wirklich noch Energie hast, Gitarre zu spielen, wow, dann alle Achtung,
ich werde dich dafür bewundern.»

Joh schubste Louis von der Seite an und schmunzelte.
Seine Bemerkung liess Louis’ Verunsicherung wachsen. Er hatte nur eine
vage Ahnung, was ihn auf dieser Alp erwartete. 
Schon Johs detaillierte Beschreibung seiner neuen Aufgaben hatte ihn
zunehmend aus der Fassung gebracht. Ob er der Arbeit eines Schäfers,
oder Zusenns, oder was immer er sein würde, gewachsen war?
Genug. Er musste aufhören, sich weitere Gedanken zu machen. Der
Arbeitsvertrag war unterschrieben. Er stand am Anfang eines neuen Weges. 
Abzweigungen gab es keine.
  
\sectionbreak

Mit dem Eintreten ins \emph{Schäferloft} knickte Louis ein. Die zurückgelassene
Verlorenheit schien lauernd hinter der Tür gesessen zu haben - nein, jetzt nicht. Er griff entschlossen nach seiner Gitarre. Die Zeit bis um sechs gehörte ihm. Manuela sei bald zurück, hatte Joh erklärt. Sie würden gemeinsam zu Abend essen. Und Joh erinnerte Louis an sein Versprechen, ihm und Manuela etwas vorzuspielen.

«Einen Musiker als Schäfer hatten wir noch nie auf der \emph{Sunnewaid},
das ist eine Première, Ciro. Ich freue mich und auch auf einen
gemütlichen Tagesabschluss zu dritt.»

Joh hatte überschwänglich geklungen.

Ein Bein angewinkelt sass Louis auf der Bettkante. Er rieb sich die Augen und senkte den Kopf. Bitte, bitte, lasst Joh keine Zeit finden, meine Anstellung zu bearbeiten. 

Louis durchstöberte nachdenklich eine Auswahl an geeigneten Songs. Johs hatte von \emph{enttäuschten Liebenden}
gesprochen\emph{, was jeder kenne.} Und Louis hatte ihnen vom Ende
seiner Beziehung erzählt. 
Beim Stimmen der Gitarre kam die Entscheidung. Passengers \emph{«Let her go»}, das
war’s.
Er mochte den Song. Hatte seine eigene Version gefunden. Immer mal
wieder hatte er ihn im Ristorante vorgespielt. Seinen Leuten und 
gelegentlich auch ausgesuchten Gästen. Es war lange her. Bevor er das Lied den beiden präsentierte, wollte er es für sich spielen. Der Text musste sitzen. Wie er sich die einzelnen Sätze in Erinnerung rief, 
stellte er sich vor, wie es gekommen wäre, wenn er Giorgia zur passenden Zeit hätte gehen lassen oder, wenn
\emph{er} der gewesen wäre, der \emph{sie} verlassen hätte. Wie ihre
Leben sich entwickelt hätten. Stark wäre er aus dieser Beziehung rausgestiegen. Länger schon hatte er geahnt, dass ihre Wege sich trennen würden. Es war hart, es sich einzugestehen. An einem
ungewissen Punkt waren sie zusammen endgültig stecken geblieben. Oder nur er? Die
Grösse jedenfalls, die er sich eben ausmalte, hätte er nicht annähernd aufgebracht. 

Er legte die Gitarre neben sich aufs Bett, streckte sich und seufzte. Seine innere Leere war im Grunde nicht ausschliesslich durch Giorgias Abkehr entstanden. Der Gedanke ploppte ganz selbstverständlich an die
Oberfläche. Es war weit gravierender. 

Sie hatte auch gleichzeitig seine Würde mitgenommen.
